% Part: normal-modal-logic
% Chapter: frame-correspondence
% Section: first-order-definability

\documentclass[../../../include/open-logic-section]{subfiles}

\begin{document}

\olfileid{nml}{frd}{fol}

\olsection{প্রথম-ক্রমে সংজ্ঞায়নযোগ্যতা}

দেখেছি, কাঠামোর অভিগম্যতা সম্পর্কের কয়েকটি ধর্ম মোডাল !!{formula}s দিয়ে সংজ্ঞায়িত করা যায়। যেমন, কাঠামোর প্রতিসাম্য !!{formula}~\Ax{B}, অর্থাৎ $p \lif \Box\Diamond
p$, দিয়ে সংজ্ঞায়িত করা যায়। এ পর্যন্ত পাওয়া সব শর্তই একটিমাত্র দ্বিস্থানের !!{predicate}-বিশিষ্ট !!a{language}-এর প্রথম-ক্রমের !!{formula}s দিয়ে প্রকাশ করা যায়। যেমন, প্রতিসাম্য $\lforall[x][\lforall[y][(\Atom{Q}{x,y} \lif \Atom{Q}{y, x})]]$ দিয়ে সংজ্ঞায়িত হয়—এই অর্থে যে, $\Domain{M} =
W$ ও $\Assign{Q}{M} = R$-বিশিষ্ট প্রথম-ক্রমের !!{structure}~$\Struct{M}$ আগের সূত্রটি পূরণ করে তখনই এবং কেবল তখনই, যখন $R$ প্রতিসম। এ থেকে পরের সংজ্ঞাটি আসে:

\begin{defn}
  কাঠামোর শ্রেণি~$\mClass{F}$ \emph{প্রথম-ক্রমে সংজ্ঞায়নযোগ্য}, যদি একটিমাত্র দ্বিস্থানের !!{predicate}~$Q$-বিশিষ্ট প্রথম-ক্রমের ভাষায় এমন একটি !!a{sentence}~$!A$ থাকে যে, $\mModel{F} =
  \tuple{W, R} \in \mClass{F}$ তখনই এবং কেবল তখনই, যখন $\Domain{M} = W$ ও $\Assign{Q}{M}
  = R$-বিশিষ্ট প্রথম-ক্রমের !!{structure}~$\Struct{M}$-এ $\Sat{M}{!A}$।
\end{defn}

এ পর্যন্ত আলোচিত ধর্ম ও তাদের সংজ্ঞায়ক মোডাল !!{formula}s বিশেষ ঘটনা। প্রত্যেক !!{formula} প্রথম-ক্রমে সংজ্ঞায়নযোগ্য কাঠামোর শ্রেণি সংজ্ঞায়িত করে না; আবার প্রথম-ক্রমে সংজ্ঞায়নযোগ্য প্রত্যেক কাঠামোর শ্রেণিও কোনো মোডাল সূত্র দিয়ে সংজ্ঞায়িত হয় না।

প্রথম দাবির বিপরীত-উদাহরণ লোবের সূত্র:
\begin{equation}
\Box(\Box p \lif p) \lif \Box p. \tag{\Ax{W}}
\end{equation}
\Ax{W} পরিযায়ী ও বিপরীতক্রমে সুপ্রতিষ্ঠিত কাঠামোর শ্রেণি সংজ্ঞায়িত করে। কোনো সম্পর্কে এমন কোনো অসীম ধারা $w_1$, $w_2$, \dots{} না থাকলে সেটি সুপ্রতিষ্ঠিত, যার জন্য $Rw_2w_1$, $Rw_3w_2$, \dots। যেমন,~$\Nat$-এর উপর $<$ সম্পর্ক সুপ্রতিষ্ঠিত, কিন্তু~$\Int$-এর উপর $<$ সম্পর্কটি তা নয়। কোনো সম্পর্কের বিপরীত সম্পর্ক সুপ্রতিষ্ঠিত হলে মূল সম্পর্কটি বিপরীতক্রমে সুপ্রতিষ্ঠিত। তাই বিপরীতক্রমে সুপ্রতিষ্ঠিত সম্পর্কে এমন কোনো অসীম ধারা $w_1$, $w_2$, \dots{} নেই, যার জন্য $Rw_1w_2$, $Rw_2w_3$, \dots।

কিন্তু পরিযায়ী ও বিপরীতক্রমে সুপ্রতিষ্ঠিত সম্পর্ক সংজ্ঞায়িত করে এমন কোনো প্রথম-ক্রমের !!{formula} নেই। বিপরীতটি ধরে নিই:~$\Sat{M}{!F}$ তখনই এবং কেবল তখনই, যখন $R =
\Assign{Q}{M}$ পরিযায়ী ও বিপরীতক্রমে সুপ্রতিষ্ঠিত। $n\ge 2$-এর জন্য $!A_n$ হলো !!{formula}
\[
(\Atom{Q}{a_1,a_2} \land \dots \land
    \Atom{Q}{a_{n-1},a_{n}})
\]
এবার !!{formula}s-এর সমষ্টি ধরি
\[
\Gamma = \{!F, !A_2, !A_3, \dots\}.
\]
$\Gamma$-র প্রত্যেক সসীম উপসমষ্টি সন্তোষণযোগ্য: উপসমষ্টিতে উপস্থিত $!A_k$-গুলির মধ্যে বৃহত্তম সূচক $k$ নিই (এমন সূত্র না থাকলে $k=2$ ধরি), $\Domain{M_k} = \{1, \dots, k\}$, $i\le k$-এর জন্য $\Assign{a_i}{M_k} = i$ এবং অন্য ধ্রুবকগুলির মান নির্বিচারে নিই; আর $\Assign{Q}{M_k} = <$। $\{1,
\dots, k\}$-এর উপর $<$ পরিযায়ী ও বিপরীতক্রমে সুপ্রতিষ্ঠিত বলে $\Sat{M_k}{!F}$। নির্মাণ অনুযায়ী, সব $i \le k$-এর জন্য $\Sat{M_k}{!A_i}$, যেখানে $i\ge2$। প্রথম-ক্রমের যুক্তির সংহতি উপপাদ্য অনুযায়ী, কোনো !!{structure}~$\Struct{M}$-এ $\Gamma$ সন্তোষণযোগ্য। অনুমান অনুযায়ী, $\Sat{M}{!F}$ বলে $\Assign{Q}{M}$ বিপরীতক্রমে সুপ্রতিষ্ঠিত। কিন্তু $\Assign{a_1}{M}$, $\Assign{a_2}{M}$, \dots{} স্পষ্টতই সেই নিষিদ্ধ ধরনের একটি অসীম ধারা গঠন করবে।
% BN-SRC-345: অনির্ধারিত সূচনা-দফা সরিয়ে শৃঙ্খল-সূত্রের সূচক দুই থেকে শুরু করা হয়েছে।

দ্বিতীয় দাবির বিপরীত-উদাহরণ সর্বজনীনতার ধর্ম: সব $u$ ও $v$-এর জন্য $Ruv$। সর্বজনীন কাঠামো প্রথম-ক্রমের সূত্র $\lforall[x][\lforall[y][\Atom{Q}{x,y}]]$ দিয়ে সংজ্ঞায়িত হয়। কিন্তু সব এবং কেবল সর্বজনীন কাঠামোতে বৈধ এমন কোনো মোডাল সূত্র নেই। এটি একটি স্বতন্ত্রভাবে গুরুত্বপূর্ণ ফলের অনুবর্তী: সর্বজনীন কাঠামোতে বৈধ সূত্রগুলি ঠিক সেই সূত্রগুলিই, যেগুলি প্রতিবিম্বী, প্রতিসম ও পরিযায়ী কাঠামোতে বৈধ। এমন প্রতিবিম্বী, প্রতিসম ও পরিযায়ী কাঠামো আছে, যেগুলি সর্বজনীন নয়। তাই সব সর্বজনীন কাঠামোতে বৈধ প্রত্যেক !!{formula} কোনো-না-কোনো অ-সর্বজনীন কাঠামোতেও বৈধ।

\end{document}
